\subsection{Scenario Definition}
We apply the framework to the motivating example from Section 1. The zones are defined as \texttt{LivingRoom}, \texttt{Hallway}, \texttt{StairEntrance}, and \texttt{Stairs}. The initial state is the \texttt{LivingRoom}, and the hazard is \texttt{Stairs}.

\subsection{Unprotected vs Protected Scenario}
In the unprotected scenario, all edges are unguarded, and the hazard is reachable. We define a \texttt{gateIntervention} that restricts the transition from \texttt{StairEntrance} to \texttt{Stairs} under the following assumptions:
\begin{itemize}
    \item The gate is installed.
    \item The gate is latched.
    \item The child cannot operate the latch.
    \item The child cannot climb over the gate.
    \item The barrier has not failed.
    \item There is no alternate route.
\end{itemize}

\subsection{Formal Verification}
We formalise this intervention and prove its soundness in Lean 4:
\begin{lstlisting}[language=Lean]
theorem stair_gate_safety (gamma : StairAssumptions)
    (h_assumptions : gamma.gateInstalled ∧ gamma.gateLatched ∧ gamma.childCannotOperateLatch ∧ gamma.childCannotClimbOver ∧ gamma.barrierNotFailed ∧ gamma.noAlternateRoute)
    (h_reach_soundness : ∀ s₀ s, s₀.zone ∈ stairInitial → Reachable (RestrictedScenario UnprotectedStaircase gateIntervention) gamma s₀ s → ¬IsHazard (RestrictedScenario UnprotectedStaircase gateIntervention) s) :
    SoundIntervention UnprotectedStaircase gateIntervention gamma
\end{lstlisting}
This corollary demonstrates how specific physical interventions translate directly into formal theorems.
